% ---------------------------------------------------------------------
\subsubsection{Requisiti Non Funzionali}
\label{subsec:requisiti-non-funzionali}
% ---------------------------------------------------------------------

Questa sottosezione descrive i requisiti non funzionali del sistema \productname{} secondo lo standard IEEE 830-1998 e il modello di qualità ISO/IEC 25002. I requisiti non funzionali specificano attributi di qualità del sistema, vincoli sulle modalità implementative e proprietà che il sistema deve soddisfare indipendentemente dalle sue funzionalità specifiche.

\vspace{0.3cm}

I requisiti non funzionali sono spesso più critici dei requisiti funzionali: mentre gli utenti possono trovare soluzioni alternative a una funzionalità sub-ottimale, il mancato rispetto di un requisito non funzionale può rendere il sistema completamente inutilizzabile.

\vspace{0.3cm}

I requisiti sono organizzati secondo la classificazione:

\begin{itemize}[leftmargin=1.1em]
\item \textbf{Product Requirements}: requisiti che specificano caratteristiche intrinseche del prodotto software (performance, usabilità, affidabilità, sicurezza)
\item \textbf{Organizational Requirements}: requisiti derivanti da politiche e procedure dell'organizzazione (linguaggi obbligatori, standard da rispettare, processi operativi)
\item \textbf{External Requirements}: requisiti derivanti da fonti esterne (conformità normativa, requisiti etici, requisiti legali)
\end{itemize}

% =====================================================================
% PRODUCT REQUIREMENTS
% =====================================================================

\paragraph{Product Requirements}
\label{subsubsec:product-requirements}

I Product Requirements definiscono le caratteristiche di qualità intrinseche del sistema \productname, basate sul modello ISO/IEC 25002.

\vspace{0.5cm}

\newpage
% =====================================================================
% REQ-NFR-01: PERFORMANCE
% =====================================================================

\begin{reqbox}[title=REQ-NFR-01: Performance]
\label{subsubsec:req-nfr-performance}

\subparagraph{Descrizione}
Il sistema deve garantire tempi di risposta rapidi per le operazioni più comuni.

\subparagraph{Requisiti}
\begin{itemize}[leftmargin=1.1em]
\item Le operazioni di visualizzazione (dashboard, lista issue, dettagli) devono completarsi entro \textbf{1 secondo}
\item Le operazioni di creazione e modifica (nuova issue, commento) devono completarsi entro \textbf{2 secondi}
\item La generazione di report mensili deve completarsi entro \textbf{5 secondi}
\item Il caricamento iniziale dell'applicazione deve completarsi entro \textbf{3 secondi}
\end{itemize}

\end{reqbox}

\vspace{0.5cm}

% =====================================================================
% REQ-NFR-02: USABILITÀ
% =====================================================================

\begin{reqbox}[title=REQ-NFR-02: Usabilità]
\label{subsubsec:req-nfr-usability}

\subparagraph{Descrizione}
Il sistema deve essere facile da usare e intuitivo per ridurre il tempo di apprendimento.

\subparagraph{Requisiti}
\begin{itemize}[leftmargin=1.1em]
\item L'interfaccia deve essere chiara e auto-esplicativa senza necessità di manuale utente
\item Le operazioni più comuni (crea issue, visualizza lista, commenta) devono essere accessibili in massimo \textbf{3 click}
\item I form devono mostrare messaggi di errore chiari quando l'input non è valido
\item Le azioni distruttive (elimina progetto, elimina utente) devono richiedere conferma esplicita
\end{itemize}

\end{reqbox}

\vspace{0.5cm}

\newpage

% =====================================================================
% REQ-NFR-03: AFFIDABILITÀ
% =====================================================================

\begin{reqbox}[title=REQ-NFR-03: Affidabilità]
\label{subsubsec:req-nfr-reliability}

\paragraph{Descrizione}
Il sistema deve funzionare in modo stabile senza crash o perdita di dati.

\paragraph{Requisiti}
\begin{itemize}[leftmargin=1.1em]
\item Il sistema non deve presentare crash durante l'uso normale
\item Eventuali errori del server devono essere gestiti mostrando messaggi chiari all'utente invece di bloccare l'applicazione
\item I dati inseriti dall'utente non devono mai essere persi (ogni operazione deve essere persistita nel database)
\end{itemize}

\end{reqbox}

\vspace{0.5cm}

% =====================================================================
% REQ-NFR-04: SICUREZZA
% =====================================================================

\begin{reqbox}[title=REQ-NFR-04: Sicurezza]
\label{subsubsec:req-nfr-security}

\paragraph{Descrizione}
Il sistema deve proteggere i dati degli utenti e impedire accessi non autorizzati.

\paragraph{Requisiti}
\begin{itemize}[leftmargin=1.1em]
\item Le password devono essere salvate come hash, mai in chiaro
\item Le sessioni devono essere gestite con durata limitata
\item Ogni richiesta API deve verificare che l'utente sia autenticato e autorizzato
\item Gli utenti possono accedere solo ai progetti di cui fanno parte
\item Solo gli amministratori possono creare progetti, gestire utenti e assegnare issue
\item La comunicazione in produzione deve avvenire su HTTPS
\item Gli input utente devono essere validati per prevenire SQL Injection
\end{itemize}

\end{reqbox}

\vspace{0.5cm}

\newpage

% =====================================================================
% REQ-NFR-05: MANUTENIBILITÀ
% =====================================================================

\begin{reqbox}[title=REQ-NFR-05: Manutenibilità]
\label{subsubsec:req-nfr-maintainability}

\paragraph{Descrizione}
Il codice deve essere organizzato in modo chiaro per facilitare manutenzione e modifiche future.

\paragraph{Requisiti}
\begin{itemize}[leftmargin=1.1em]
\item Il backend deve essere organizzato in layer separati: Controller, Service
\item Il frontend deve utilizzare componenti React riusabili e ben organizzati
\item Il codice deve seguire convenzioni di naming consistenti
\item Devono essere presenti commenti per spiegare logiche complesse
\item Il progetto deve avere un README con istruzioni per setup e deployment
\item Il codice deve essere formattato automaticamente (Prettier per TypeScript, formatter Java)
\end{itemize}

\end{reqbox}

\vspace{0.5cm}

\newpage

% =====================================================================
% REQ-NFR-06: COMPATIBILITÀ
% =====================================================================

\begin{reqbox}[title=REQ-NFR-06: Compatibilità]
\label{subsubsec:req-nfr-compatibility}

\paragraph{Descrizione}
Il sistema deve funzionare correttamente sui browser più comuni.

\paragraph{Requisiti}
\begin{itemize}[leftmargin=1.1em]
\item L'applicazione deve funzionare su Chrome, Firefox, Edge e Safari versioni recenti
\item L'interfaccia deve essere responsive e utilizzabile su qualsiasi schermo da 320px in su
\item Le API REST devono restituire JSON strutturato con status code HTTP corretti
\end{itemize}

\end{reqbox}

\vspace{0.5cm}

% =====================================================================
% REQ-NFR-07: CONFORMITÀ NORMATIVA
% =====================================================================

\begin{reqbox}[title=REQ-NFR-07: Conformità Normativa]
\label{subsubsec:req-nfr-legal}

\paragraph{Descrizione}
Il sistema deve rispettare le normative di base sulla protezione dei dati.

\paragraph{Requisiti}
\begin{itemize}[leftmargin=1.1em]
\item Il sistema deve raccogliere solo i dati necessari (nome, cognome, email, password)
\item I dati personali devono essere protetti e accessibili solo agli utenti autorizzati
\item Gli amministratori possono eliminare account utente quando richiesto
\item Le dipendenze open-source devono avere licenze compatibili (MIT, Apache 2.0)
\end{itemize}

\end{reqbox}

\vspace{0.5cm}

\newpage

% =====================================================================
% RIEPILOGO
% =====================================================================

\subsubsection{Riepilogo}
\label{subsubsec:riepilogo-nfr}

La seguente tabella riassume i requisiti non funzionali con priorità:

\begin{longtable}{|p{0.20\linewidth}|p{0.50\linewidth}|p{0.20\linewidth}|}
\hline
\textbf{ID} & \textbf{Requisito} & \textbf{Priorità} \\ \hline
\endfirsthead

\hline
\textbf{ID} & \textbf{Requisito} & \textbf{Priorità} \\ \hline
\endhead

\hline
\endfoot

REQ-NFR-01 & Performance & Alta \\ \hline
REQ-NFR-02 & Usabilità & Alta \\ \hline
REQ-NFR-03 & Affidabilità & Media \\ \hline
REQ-NFR-04 & Sicurezza & Critica \\ \hline
REQ-NFR-05 & Manutenibilità & Media \\ \hline
REQ-NFR-06 & Compatibilità & Media \\ \hline
REQ-NFR-07 & Conformità Normativa & Alta \\ \hline

\end{longtable}